
\section{Canonical integration by parts and Gamma parity}
\label{ibp:section}

The pole calculus determines a finite part from a differential polynomial;
it does not yet identify a preferred polynomial representative.  The
three continuations in \cite{dpm:reports} also give compatible answers
to that formal question.  We retain one triangular basis, its explicit
endpoint functional, and the inverse-Gamma parity argument.  This is a
quotient of vector spaces before evaluation in special values.

Let $\mathcal A=\mathbb Q[u_0,u_1,\ldots]$, with each polynomial
using finitely many variables, and set
\[
 \mathcal D=\sum_{r\ge0}u_{r+1}\frac{\partial}{\partial u_r},
 \qquad \operatorname{wt}(u_r)=r+1.
\]
The factor degree of a monomial is its number of factors, counted with
multiplicity.  The derivation increases weight by one and preserves
factor degree.  Under $u_r=\psi^{(r)}(x)$ it becomes $d/dx$.
Write $\mathcal A_w$ for weight $w$ and $p(w)$ for the partition
number, so $\dim\mathcal A_w=p(w)$.
For $w\ge2$, let $\mathcal B_w$ be the monomials whose largest
variable index occurs at least twice; this includes $u_0^w$.
At weights zero and one use $\mathcal B_0=\{1\}$ and
$\mathcal B_1=\{u_0\}$ separately.

\begin{theorem}[A unique formal normal form at every weight]
\label{ibp:normal}
For $w\ge2$,
\begin{equation}\label{ibp:direct}
 \mathcal A_w=\mathcal D\mathcal A_{w-1}
          \oplus\operatorname{span}_{\mathbb Q}\mathcal B_w.
\end{equation}
Thus every $P\in\mathcal A_w$ has a unique decomposition
$P=\mathcal D Q+N$, with $Q\in\mathcal A_{w-1}$ and $N$ in
the displayed normal space.  A terminating rational algorithm computes
both terms.  If $p_\ell(w)$ counts partitions of $w$ with exactly
$\ell$ parts, the dimension at factor degree $\ell$ is
$p_\ell(w)-p_\ell(w-1)$.  In particular
\[
 |\mathcal B_w|=p(w)-p(w-1),\qquad
 \sum_{w\ge0}\dim(\mathcal A_w/\mathcal D\mathcal A_{w-1})t^w
       =t+\prod_{j\ge2}(1-t^j)^{-1}.
\]
Every retained positive derivative has order at most
$\lfloor w/2\rfloor-1$.
\end{theorem}

\begin{proof}
A monomial outside the normal space has a unique largest index
$r\ge1$.  Write it as $u_r u_{r-1}^a V$, where $V$ uses indices
at most $r-2$.  The product rule gives
\begin{equation}\label{ibp:step}
 u_r u_{r-1}^a V
 =\mathcal D\left(\frac{u_{r-1}^{a+1}V}{a+1}\right)
       -\frac{u_{r-1}^{a+1}\mathcal D V}{a+1}.
\end{equation}
The remainder has strictly smaller largest index, and factor degree
is unchanged.  Repeated application terminates in the finite set of
partitions of $w$, proving spanning and the algorithm.

For directness, a nonconstant $Q$ with largest index $m$ has in
$\mathcal DQ$ the nonzero terms
$u_{m+1}\partial Q/\partial u_m$.  These have largest index $m+1$
exactly once; the other terms have smaller largest index.  No element
of the normal space has this support.  Characteristic zero gives
$\partial Q/\partial u_m\ne0$, so the intersection is zero and
$\mathcal D$ is injective on positive weights.  Equivalently, increasing
one largest part of a partition of $w-1$ gives a distinct leading
partition with unique largest part at weight $w$.  Decreasing that
unique part is its inverse.  The bijection preserves the number of
parts, proving both dimension formulas.  The constant kernel at
weight zero supplies the exceptional $t$ in the generating series.
A repeated largest part has size at most $w/2$, giving the derivative
order bound.
\end{proof}

For instance the bases at weights four through seven are
\[
\begin{array}{c|l}
4&u_0^4,\ u_1^2\\
5&u_0^5,\ u_0u_1^2\\
6&u_0^6,\ u_0^2u_1^2,\ u_1^3,\ u_2^2\\
7&u_0^7,\ u_0^3u_1^2,\ u_0u_1^3,\ u_0u_2^2.
\end{array}
\]
The image $\mathcal D\mathcal A$ is not generally an ideal.
For example $u_1=\mathcal D u_0$, whereas $u_1^2$ is a nonzero
normal representative at weight four.  Multiplication of quotient
classes is therefore not being defined.  The dimensions count formal
representatives, without asserting independence after integration.

\subsection{The endpoint functional is retained}

For a differential polynomial $Q$, write $f_Q$ for its digamma
substitution as in Section~\ref{dpm:section}, and define
\begin{equation}\label{ibp:endpoint}
 \mathcal E(Q)=f_Q(1)-[x^0]f_Q(x).
\end{equation}
The coefficient includes products between pole and regular terms.
It is not the product of the individual constant coefficients.

\begin{corollary}[Exact finite-part integration by parts]
\label{ibp:integral}
For $P=\mathcal DQ+N$ from Theorem~\ref{ibp:normal},
\begin{equation}\label{ibp:FP}
 \operatorname{FP}_x\int_0^1 f_P(x)\,dx
       =\mathcal E(Q)+\operatorname{FP}_x\int_0^1f_N(x)\,dx.
\end{equation}
For $\operatorname{Re}a>0$, the corresponding ordinary shifted
identity retains both endpoints:
\[
 \int_0^1 f_P(a+x)\,dx
  =f_Q(a+1)-f_Q(a)+\int_0^1 f_N(a+x)\,dx.
\]
\end{corollary}

\begin{proof}
Integrate $f_{\mathcal DQ}=f_Q'$ over $[\varepsilon,1]$ before
taking the constant term.  The answer is $f_Q(1)-f_Q(\varepsilon)$,
whose constant is exactly \eqref{ibp:endpoint}.  Linearity of the
fixed-coordinate finite part proves \eqref{ibp:FP}.  With the stated
shift the interval avoids all poles, and the fundamental theorem of
calculus gives the ordinary identity.
\end{proof}

The functional is explicitly finite.  For homogeneous $Q$ of weight
$W>0$, substitute the differentiated truncation
\[
 -x^{-1}-\gamma+\sum_{j=2}^{W}(-1)^j\zeta(j)x^{j-1}
\]
to compute $[x^0]f_Q(x)$.  An omitted digamma term starts at $x^W$;
replacing a leading pole by it in a weight-$W$ product raises the total
exponent by at least $W+1$, so it cannot contribute a constant.
The other endpoint uses
$\psi(1)=-\gamma$ and
$\psi^{(r)}(1)=(-1)^{r+1}r!\zeta(r+1)$ for $r\ge1$.
Consequently $\mathcal E(Q)$ lies in
$\mathbb Q[\gamma,\zeta(2),\ldots,\zeta(W)]$.

Three useful exact examples are
\[
 \mathcal E(u_0^2)=2\zeta(2),\qquad
 \mathcal E(u_0u_1)=-\zeta(3),\qquad
 \mathcal E(u_1u_2)=4\zeta(5).
\]
Thus dropping a total derivative would lose a nonzero endpoint.
Together with the trigamma-square calculation in
Section~\ref{dpm:section}, the middle identity gives
\begin{equation}\label{ibp:digamma-trigamma}
 \operatorname{FP}_x\int_0^1\psi(x)\psi''(x)\,dx
          =-3\zeta(3)-4\zeta'(3).
\end{equation}
Also $[x^0]\psi'(x)^2=\zeta(2)^2+6\zeta(4)$, so integrating
$(\psi'^2)'/2$ gives
\begin{equation}\label{ibp:trigamma-tetragamma}
 \operatorname{FP}_x\int_0^1\psi'(x)\psi''(x)\,dx=-3\zeta(4).
\end{equation}
The identities are proved by finite pole algebra and endpoint extraction;
independent native quadrature is a separate diagnostic.

\subsection{Inverse-Gamma normal forms have factor parity}

Define $A_0=1$ and $A_{m+1}=u_0A_m-\mathcal DA_m$.  After
substitution these are the differential coefficients of the Gamma ratio:
\[
 A_m(x)=(-1)^m\Gamma(x)\frac{d^m}{dx^m}\Gamma(x)^{-1},\qquad
 \frac{\Gamma(x)}{\Gamma(x-t)}
       =\sum_{m\ge0}A_m(x)\frac{t^m}{m!}.
\]
Let $A_m=N_m+\mathcal DT_m$ be the unique normal decomposition,
with $N_0=1$, $N_1=u_0$, $T_0=T_1=0$.

\begin{theorem}[All-degree parity and a recurrence retaining the primitive]
\label{ibp:parity}
Every monomial in $N_m$ has factor degree congruent to $m$ modulo two.
For $m\ge1$, decompose
\[
 u_0N_m-u_1T_m=R_m+\mathcal DS_m
\]
by Theorem~\ref{ibp:normal}.  Then
\begin{equation}\label{ibp:pairrecurrence}
 N_{m+1}=R_m,\qquad
 T_{m+1}=S_m+u_0T_m-N_m-\mathcal DT_m.
\end{equation}
All operations are finite and rational.
\end{theorem}

\begin{proof}
Introduce $p=e^F$, $q=e^{-F}$ with $F'=u_0$ only as a convenient
notation: the products below are differential polynomials because the
exponential factors cancel.  Set
\[
 L_m=\sum_{j=0}^{m-1}(-1)^j p^{(j)}q^{(m-1-j)}.
\]
The product rule telescopes to
\[
 \mathcal DL_m=pq^{(m)}+(-1)^{m-1}p^{(m)}q
         =(-1)^m A_m(u)-A_m(-u).
\]
This is also a purely formal identity: replace
$p^{(j)}/p$ by $(-1)^jA_j(-u)$ and
$q^{(j)}/q$ by $(-1)^jA_j(u)$.
Projection onto the normal space annihilates the derivative and
preserves factor degree.  Therefore
$N_m(-u)=(-1)^mN_m(u)$, proving parity.
For the recurrence, substitute $A_m=N_m+\mathcal DT_m$ and use
$u_0\mathcal DT_m=\mathcal D(u_0T_m)-u_1T_m$.
The result is
\[
 A_{m+1}=u_0N_m-u_1T_m
           +\mathcal D(u_0T_m-N_m-\mathcal DT_m).
\]
Uniqueness gives \eqref{ibp:pairrecurrence}.
\end{proof}

The first informative remainders are
\begin{align*}
 N_4&=u_0^4-u_1^2,\qquad
 N_5=u_0^5-5u_0u_1^2,\\
 N_6&=u_0^6-15u_0^2u_1^2+u_2^2,\\
 N_7&=u_0^7-35u_0^3u_1^2+7u_0u_2^2,\\
 N_8&=u_0^8-70u_0^4u_1^2+28u_0^2u_2^2-7u_1^4-u_3^2.
\end{align*}
In particular $u_1^3$, although a valid weight-six normal coordinate,
is excluded from $N_6$ by parity.  Every identity
$A_m=N_m+\mathcal DT_m$ yields an exact finite-part identity with
$\mathcal E(T_m)$ retained.  No vanishing endpoint or arithmetic
independence is inferred from this formal cancellation.

Independent rational matrix certificates verify the two-sided basis
inverses and all monomial decompositions through weight fourteen.
Separate exponential Bell and Lagrange-concomitant constructions check
the inverse-Gamma coefficients through degree twelve.  They check the
implementation and the printed examples; the preceding proofs establish
the all-weight statements.  The Gamma--Appell transport is proved in Section~\ref{gat:section};
explicit sixth/seventh spectral-coordinate reductions are given in
Section~\ref{hsm:section}. Further arithmetic reduction remains open.
